\documentclass[11pt,a4paper]{article}
\usepackage{../shared/hanzocover}
\usepackage[utf8]{inputenc}
\usepackage[T1]{fontenc}
% Shared preamble for the compute-market series.
% One style, one vocabulary, inputted by every paper so the series reads as one voice.
\usepackage[margin=1in]{geometry}
\usepackage{booktabs}
\usepackage{amsmath}
\usepackage{amssymb}
\usepackage{amsthm}
\usepackage{graphicx}
\usepackage{xcolor}
\usepackage{hyperref}
\hypersetup{colorlinks=true, linkcolor=blue, urlcolor=blue, citecolor=blue}
\usepackage{enumitem}
\usepackage{listings}
\usepackage{array}
\usepackage{tabularx}
\usepackage{siunitx}
\sisetup{group-separator={,}, group-minimum-digits=4}

\definecolor{kw}{rgb}{0.13,0.29,0.53}
\definecolor{cmt}{rgb}{0.40,0.40,0.40}
\definecolor{str}{rgb}{0.16,0.50,0.16}
\lstset{
  basicstyle=\ttfamily\footnotesize,
  keywordstyle=\color{kw}\bfseries,
  commentstyle=\color{cmt}\itshape,
  stringstyle=\color{str},
  showstringspaces=false,
  breaklines=true,
  columns=fullflexible,
  frame=single,
  framesep=4pt,
}

\newtheorem{definition}{Definition}
\newtheorem{proposition}{Proposition}
\newtheorem{remark}{Remark}

\newcommand{\code}[1]{\texttt{#1}}
\newcommand{\repo}[1]{\texttt{#1}}

% Series vocabulary — used identically in every paper.
\newcommand{\job}{\ensuremath{J}}
\newcommand{\bundle}{\ensuremath{R}}
\newcommand{\cost}{\ensuremath{\mathcal{C}}}

\tolerance=800
\emergencystretch=3em


\title{Hybrid Post-Quantum Retrofit:\\Upgrading a Live Protocol Without a Fork}

\author{
Hanzo AI Research\thanks{Corresponding author: research@lux.network} \\
\textit{Lux Industries Inc} \\
\textit{A subsidiary of Hanzo Industries Inc} \\
\texttt{https://lux.network}
}

\date{2026}

\begin{document}
\hanzocoverpage

\begin{abstract}
A compute market protects its users' funds with signatures and its users' work
with encryption, and both rest on elliptic curves. Neither survives a
cryptographically relevant quantum computer. The usual objection to acting now
is that migration means a fork, a flag day, and a coordination problem across an
open provider network.

It does not have to. We describe a retrofit pattern that adds post-quantum
identity and key exchange to a running protocol without changing its wire format
or forking the reference implementation, and we report two independent
deployments of it: one onto a mesh networking stack with a 500-byte maximum
transmission unit and a 5 bit-per-second floor, and one onto a
defence-grade file storage and transfer profile. The first is the harder
constraint by a wide margin, and it fits.

We then apply the pattern surface by surface to a decentralized compute market
--- provider identity, session establishment, job receipts, staking signatures,
governance --- give the byte cost of each, and propose a migration order derived
from key lifetime rather than from implementation convenience.
\end{abstract}

\tableofcontents
\newpage

\section{Two clocks, not one}

Quantum risk is usually discussed as a single future event. For protocol design
it is two separate deadlines with different urgencies.

\paragraph{Confidentiality has already expired.} An adversary recording session
transcripts today can decrypt them whenever a capable machine exists. Every
session established with a classical key exchange is retroactively readable. The
deadline for key exchange was whenever recording began.

\paragraph{Authenticity expires later, but the key does not rotate.} A signature
forged after the fact is worthless against a transcript already accepted, so
signing is less urgent --- unless the key is long-lived, in which case an
adversary who recovers it can impersonate the holder from that moment forward.
Provider identity keys, staking keys and governance keys are exactly that: keys
whose value is their continuity.

The ordering falls out. Key exchange first because the loss is already
accruing; long-lived identity next because it cannot be rotated cheaply;
ephemeral per-job signatures last because they can simply be reissued.

\section{The pattern}

The retrofit is additive, hybrid, and wire-compatible.

\paragraph{Additive.} The post-quantum material is carried alongside the
classical material, not instead of it. An old peer ignores a field it does not
recognise; a new peer verifies both. No version negotiation, no flag day.

\paragraph{Hybrid.} Security is the maximum of the two assumptions, not the
minimum. The session key derives from both shared secrets through a
domain-separated key derivation, so the construction is at least as strong as
the stronger component. A break in the lattice assumption degrades to classical
security; a quantum adversary degrades to lattice security. This matters because
lattice cryptography is young, and a market should not bet its custody on the
newest assumption alone.

\paragraph{Wire-compatible.} The transport, framing, session state machine and
authentication flow are untouched. This is the property that makes it a retrofit
rather than a migration.

\begin{table}[h]
\centering
\begin{tabular}{lll}
\toprule
Layer & Classical & Hybrid replacement \\
\midrule
Identity signing & Ed25519 & Ed25519 + ML-DSA-65 (detached pair, both verify) \\
Key exchange & X25519 & X-Wing = X25519 + ML-KEM-768 \\
Key exchange, high value & X448 & X448 + ML-KEM-1024 \\
Bulk AEAD & AES-256-GCM & unchanged \\
Bulk AEAD, software & ChaCha20-Poly1305 & unchanged \\
Hash & SHA-256 & SHA-384 / SHA3-384 \\
Derivation & HKDF & HKDF-SHA-384, domain-separated per layer \\
\bottomrule
\end{tabular}
\caption{The substitution table. Symmetric primitives are left alone: Grover
halves the effective key length, and 256-bit keys are already sized for that.}
\end{table}

\section{Two deployments}

\subsection{Mesh networking under extreme constraint}

The first deployment extends a cryptography-based mesh networking stack designed
for LoRa, packet radio, KISS modems, serial lines and free-space optical links
--- any half-duplex carrier with at least 5 bits per second and a 500-byte
maximum transmission unit. It provides initiator-anonymous, coordination-free
multi-hop transport entirely without IP.

Adding a 1952-byte public key and a 3309-byte signature to a protocol with a
500-byte frame is the least forgiving environment we could have chosen, which is
why it was the right test. Identity signing moved to Ed25519 with ML-DSA-65
alongside; key exchange to X25519 with ML-KEM-768; session encryption and
derivation unchanged, now over the hybrid secret. The reference protocol was not
forked.

\subsection{Defence-grade storage and transfer}

The second is an envelope-encryption profile for file storage and transfer that
encrypts before content enters any storage or transfer system, so the back end
sees only ciphertext and opaque manifests --- sender, recipients, filenames and
classification labels never leave the client. Bulk encryption is AES-256-GCM in
the certified profile or ChaCha20-Poly1305 in software, hashing is SHA-384 or
SHA3-384, key wrapping is X-Wing for the general case and X448 with ML-KEM-1024
for high-value material, and the sender signature is a detached Ed25519 and
ML-DSA-65 pair.

The two deployments share no code and no threat model. What they share is the
pattern, which is the point.

\section{Policy as a separate object}

A hybrid deployment needs a way to say ``classical is no longer acceptable
here'' without that statement being scattered across every call site. We keep it
as one value and one gate.

The policy is a profile carrying one prohibition flag per cryptographic
primitive family, with a zero value meaning classical semantics and a strict
constructor setting every flag. Operations report themselves through a small
fixed set of stable identifiers covering the precompile families and polynomial
commitments. Every constrained operation calls a single refusal function once at
the top of its entry point and otherwise proceeds in its own lane.

The property worth copying is the decomplection: the profile does not know how
operations compute, and operations do not know how the profile was selected.
Verification logic and policy enforcement are not braided together, so either can
be replaced without touching the other.

A second, duller lesson is worth stating because it cost us real bugs. The
profile and policy identifier strings were independently reimplemented in Go, in
Solidity, in TypeScript and in a paper, and they drifted --- one surface used a
name the on-chain state did not recognise. Identifiers of this kind should have
exactly one source of truth, generated into every consuming language. We now
publish them as a single package with Go, Solidity and TypeScript subpackages
sharing one set of strings.

\section{Applying it to a compute market}

\begin{table}[h]
\centering
\begin{tabular}{llll}
\toprule
Surface & Key lifetime & Exposure & Priority \\
\midrule
Session establishment & seconds & harvest-now, decrypt-later & first \\
Provider identity & years & impersonation on recovery & first \\
Staking / custody keys & years & direct theft & first \\
Governance keys & years & protocol capture & second \\
Job receipts & minutes & forgery after acceptance & third \\
Client API tokens & months & rotate instead & rotate \\
\bottomrule
\end{tabular}
\caption{Migration order derived from key lifetime and loss on compromise.}
\end{table}

\subsection{Byte budget}

Adoption is a bandwidth question, and the numbers are modest at market scale.

\begin{table}[h]
\centering
\begin{tabular}{lrr}
\toprule
Object & Classical & Hybrid \\
\midrule
Identity public key & \SI{32}{\byte} & \SI{1984}{\byte} \\
Identity signature & \SI{64}{\byte} & \SI{3373}{\byte} \\
Key-exchange public key & \SI{32}{\byte} & \SI{1216}{\byte} \\
Key-exchange ciphertext & \SI{32}{\byte} & \SI{1120}{\byte} \\
\bottomrule
\end{tabular}
\caption{Wire cost of the hybrid substitution. ML-DSA-65 public keys are
\SI{1952}{\byte} and signatures \SI{3309}{\byte}; ML-KEM-768 public keys are
\SI{1184}{\byte} and ciphertexts \SI{1088}{\byte}; classical components add
\SI{32}{\byte} each.}
\end{table}

A signed job receipt grows by roughly \SI{3.3}{\kilo\byte}. At a million jobs a
day that is \SI{3.3}{\giga\byte} of additional signature traffic --- entirely
unremarkable for a system already moving model weights. Session establishment
grows by about \SI{2.3}{\kilo\byte} once per session, amortised across every
request in it.

\subsection{A conservative root of trust}

Lattice assumptions are the right default for volume: fast, small, standardised.
For the single most valuable key in the system --- the one that anchors
governance or custody --- there is an argument for a different hardness
assumption entirely. Hash-based signatures under FIPS 205 rest only on the
preimage resistance of a hash function, an assumption with a far longer track
record and no algebraic structure for a future attack to exploit. They are larger
and slower, which is irrelevant for a key used a handful of times a year.

Our threshold signing daemon carries this as a first-class option alongside the
lattice families, and also supports composing two orthogonal lattice families so
that both must verify --- a hedge against a break in one lattice problem rather
than a break in lattices generally.

\section{What adoption looks like}

The retrofit does not require adopting our transport, our consensus or our
tooling. It requires an implementation of ML-DSA-65 and ML-KEM-768, a
domain-separated derivation, an additive wire field, and a dual-verify window
during which both signatures are checked and only the classical one is required.
The window closes when the population has migrated, at which point the classical
field becomes optional and eventually ignored.

Our implementations of the primitives, the hybrid envelope profile, the mesh
retrofit, the policy gate and the shared identifier package are open source under
permissive licence, with known-answer test vectors. They are usable
independently of each other and of everything else we build.

\end{document}
